%! Author = sriadityavedantam
%! Date = 3/25/26

\section{Derivatives and the Intermediate Value Property}

\lesson{33}{Monday 27 October 2025 9:10}{Day 33}

\begin{definition}[Derivative]
    Let $f: A\to\rr$, $c\in A$.
    Then
    \[
        f\uptick(c) = \lim_{x\to c} \dfrac{f(x) - f(c)}{x - c}.
    \]
    If this limit exists, then $f$ is differentiable at $c$.
\end{definition}

\begin{problem}{
    $f(x) = \abs{x}$ is differentiable at all $c\neq 0$.
}
    \begin{align*}
        \lim_{x\to c} \dfrac{\abs{x} - \abs{c}}{x-c}
        &= \begin{cases}
               \lim_{x\to c} \dfrac{x - c}{x-c} = 1,\quad c > 0\\
               \lim_{x\to c}\dfrac{-x + c}{x-c} = -1,\quad c < 0.
        \end{cases}
    \end{align*}

    At $c = 0$, the one-sided limits are not equal.
\end{problem}

\begin{proposition}{
    If $f$ is differentiable at $c$, then $f$ is continuous at $c$.
}
    We have
    \[
        \dfrac{f(x) - f(c)}{x-c}\to f\uptick(c)
    \]
    as $x\to c$.
    Also
    \[
        x-c\to 0.
    \]
    Hence
    \[
        f(x) - f(c) = \dfrac{f(x)-f(c)}{x-c}(x-c)\to f\uptick(c)\cdot 0 = 0.
    \]
    So
    \[
        f(x)\to f(c).
    \]
    Therefore $f$ is continuous at $c$.
\end{proposition}

\begin{theorem*)}{
    Derivative arithmetic is well defined with the usual conventions.
}
\end{theorem*}

\begin{theorem}[Linear Approximation Formula]{
    Suppose $f : A\to\rr$ and $c\in A$.
    Then $f$ is differentiable at $c$ with derivative $L$ if and only if
    \[
        f(x) = f(c) + L(x-c) + \e(x),
    \]
    where
    \[
        \lim_{x\to c} \dfrac{\e(x)}{x-c} = 0.
    \]
}
    \[
        f(x) = f(c) + L(x-c) + \e(x)
    \]
    if and only if
    \[
        \e(x) = f(x) - f(c) - L(x-c).
    \]
    Hence
    \[
        \dfrac{\e(x)}{x-c} = \dfrac{f(x) - f(c)}{x-c} - L.
    \]
    So
    \[
        \lim_{x\to c} \dfrac{\e(x)}{x-c} = 0
    \]
    if and only if
    \[
        \lim_{x\to c}\dfrac{f(x)-f(c)}{x-c} = L.
    \]
    Thus $f\uptick(c) = L$.
\end{theorem}

\begin{corollary}
    $f$ is differentiable at $c$ with derivative $L$ if and only if
    \[
        f(x) = f(c) + d(x)(x-c),
    \]
    where
    \[
        \lim_{x\to c} d(x) = L.
    \]
\end{corollary}

\begin{theorem}[Chain Rule]{
    Suppose $g:A\to\rr$ and $f:B\to \rr$, where $g(A)\subset B$.
    Suppose $g$ is differentiable at $c$ and $f$ is differentiable at $g(c)$.
    Then $f\circ g$ is differentiable at $c$ and
    \[
        (f\circ g)\uptick(c) = f\uptick(g(c))g\uptick(c).
    \]
}
    We will use the corollary.
    Write
    \[
        g(x) = g(c) + d(x)(x-c),
    \]
    where
    \[
        d(x)\to g\uptick(c)
    \]
    as $x\to c$.
    Also write
    \[
        f(y) = f(g(c)) + \delta(y)(y-g(c)),
    \]
    where
    \[
        \delta(y)\to f\uptick(g(c))
    \]
    as $y\to g(c)$.
    Substituting $y=g(x)$ gives
    \begin{align*}
        f(g(x))
        &= f(g(c)) + \delta(g(x))(g(x)-g(c))\\
        &= f(g(c)) + \delta(g(x))d(x)(x-c).
    \end{align*}
    Since $g$ is continuous at $c$, we have $g(x)\to g(c)$.
    Hence
    \[
        \delta(g(x))\to f\uptick(g(c)).
    \]
    Also
    \[
        d(x)\to g\uptick(c).
    \]
    Therefore
    \[
        \delta(g(x))d(x)\to f\uptick(g(c))g\uptick(c).
    \]
    So by the corollary,
    \[
        (f\circ g)\uptick(c) = f\uptick(g(c))g\uptick(c).
    \]
\end{theorem}

\section{Mean Value Theorems}

\begin{theorem}[Mean Value Theorem]
    Suppose $f$ is continuous on $[a,b]$ and differentiable on $(a,b)$.
    Then there exists $c\in(a,b)$ such that
    \[
        f(b) - f(a) = f\uptick(c)(b-a).
    \]
\end{theorem}

\begin{theorem}[Cauchy MVT]
    Suppose $f,g$ are continuous on $[a,b]$ and differentiable on $(a,b)$.
    Then there exists $c\in(a,b)$ such that
    \[
        f\uptick(c)(g(b) - g(a)) = g\uptick(c)(f(b) - f(a)).
    \]
\end{theorem}

\begin{theorem}[Rolle's Theorem]{
    Suppose $f$ is continuous on $[a,b]$ and differentiable on $(a,b)$, and $f(b) = f(a)$.
    Then there exists $c\in (a,b)$ such that $f\uptick(c) = 0$.
}
    We claim there exists $c\in (a,b)$ where $f$ attains either its maximum or its minimum.
    Since $f$ is continuous and $[a,b]$ is compact, the maximum and minimum are attained.
    If both are attained only at the endpoints, then since $f(a)=f(b)$, the function is constant.
    Otherwise, $f$ attains either a maximum or a minimum at some interior point $c\in(a,b)$.
    Say $f$ achieves its maximum at $c$.
    Then
    \[
        \lim_{x\to c^+}\dfrac{f(x) - f(c)}{x-c}\leq 0
    \]
    and
    \[
        \lim_{x\to c^-}\dfrac{f(x) - f(c)}{x-c}\geq 0.
    \]
    Since $f$ is differentiable at $c$, these one-sided limits are equal.
    Hence
    \[
        f\uptick(c)=0.
    \]
\end{theorem}